\chapter{Simulation Experiments}
\label{chp:simulation}

\section{Experimental Setup}

The simulation experiments are designed to probe the central challenges of single-light navigation: navigating around obstacles, and recovering a reliable gradient from the sparse measurements gathered along a flight path, which first requires that the light field forms a navigable gradient at all.

All experiments share a single setup in the Webots environment described in Chapter~\ref{chp:background}: a single overhead light source with a known position serves as the navigation target, providing a controlled setting in which each challenge can be isolated and tested (Figure~\ref{fig:sim-setup}).

\begin{figure}[t]
    \centering
    \missingfigure[figwidth=0.6\textwidth]{Webots simulation setup --- single overhead light source, drone start pose, and obstacle placement}
    \caption{Shared simulation experiment setup used across all simulation experiments.}
    \label{fig:sim-setup}
\end{figure}

\section{Experiment 1: Baseline Bearing-Only Navigation}

\textit{Question: how does the bearing-only controller behave in open space and in the presence of obstacles?}

The bearing-only controller navigates successfully toward the light source in unobstructed conditions but fails to navigate around obstacles, as it has no spatial memory of the light environment and cannot recover a path once line-of-sight is disrupted.

\section{Experiment 2: Dense Probe Map Collection}

\textit{Question: does the light intensity field form a navigable gradient, and what does it look like under ideal sampling conditions?}

A systematic grid scan produces a dense probe map of the light intensity field, confirming a consistent gradient toward the source and providing an ideal-coverage reference dataset for evaluating gradient estimation methods.

\section{Experiment 3: Gradient Estimation Method Comparison}

\textit{Question: which of the five gradient estimation methods most accurately recovers the gradient direction from the dense probe map?}

The five candidate gradient estimation methods introduced in Chapter~\ref{chp:algorithm} are evaluated against the dense probe map, with weighted least squares yielding the lowest directional error.

\textit{Question: how does gradient estimation quality degrade when the drone can only use the map points visited along its flight path?}

Gradient quality is re-evaluated using only the flight-path-visited subset of map points, quantifying the accuracy penalty of the realistic sparse coverage scenario compared to the ideal full-map case.

\section{Experiment 4: Fused Controller Validation}

\textit{Question: can the fused controller navigate around obstacles that the bearing-only controller cannot handle?}

The fused controller is validated in simulation by demonstrating successful navigation around an obstacle, with the state machine correctly transitioning through modes as the drone's map coverage and bearing signal quality evolve.
